\documentclass[11pt]{article}

\usepackage{maiacustom}

\begin{document}

\psettitle{Astronomy Question Bank}

These questions were carefully written or selected to prepare students for the astronomy team selection process in Brazil. Some questions are not original; those are marked with () before the statement. The question-bank template is the same one used by Professor Kevin Zhou. His work is invaluable, and many ideas in this set can be found in his handouts.

% \begin{psidea}{Idea title}{}
% Idea
% \end{psidea}

% \begin{psexample}{Example title}{}
% Example
% \end{psexample}

% \begin{pssolution*}{}{}
% Solution
% \end{pssolution*}

% \begin{psremark*}{Remark title}{}
% Remark
% \end{psremark*}
\section{Photometry and Modern Physics}

\pts{3} 
\begin{pproblem}
    The goal of this question is to derive an expression for the \textit{optical depth}. Imagine that a beam of light passes through a region of space containing a certain number of particles. If \(I_0\) is the intensity of the light before it passes through that region of space and \(I\) is the intensity afterward, the optical depth is defined as:
    \[I = I_0e^{-\tau}\]
    where \(\tau\) is the optical depth.
    \\
    Consider a region of space with particle number density \(n\).
    \begin{alternativas}
        \item What is the number of particles in an area \(A\) and thickness \(dz\)?
        \item Assuming that each particle has cross section \(\sigma\), what is the area blocked by the particles?
        \item Find the formula for \(\tau\). 
    \end{alternativas}

\end{pproblem}

\pts{2}
\begin{pproblem} The Triangulum Galaxy, \(M33\), is the third-largest galaxy in the Local Group. It lies at a distance \(d=970 \text{ kpc}\) from us and has an apparent magnitude of \(5.72\). Given that it contains approximately \(40\) billion stars, find the average luminosity of the stars in \(M33\). Does your estimate seem consistent with reality? Why?
\end{pproblem}

\pts{4}
\begin{pproblem}
    Consider a universe with a constant stellar number density, that is, a number of stars per unit volume, of value \(n\). Assume that every star has \(L = L_\odot\) and that the universe contains a total of \(N_0\) stars.
    
    \begin{alternativas}
    \item What is the probability that the absolute magnitude of a star, seen from the center of the universe, lies between \(m\) and \(m+dm\), where \(dm\) is an infinitesimal interval of magnitude? Leave your answer in terms of \(m\) and the largest possible magnitude, \(m_{lim}\), of a star at the "edge" of the universe.

    \item What is the probability that a star can be observed with the naked eye?
    
    \textbf{Data:} \[\int_{-\infty}^{6}10^{0.6(x-A)}dx \approx 2881.6 e^{-1.381 A}\]
    \end{alternativas}
\end{pproblem}

\pts{4}
\begin{pproblem}
    The \textit{surface brightness}, flux per solid angle per frequency, \(B_\nu\), is given by Planck's law:
    \[B_\nu = \frac{dF}{d\nu d\Omega} = \frac{2h\nu^3}{c^2(e^{\frac{h\nu}{k_BT}}-1)}\] 
    \begin{alternativas}
        \item Find an expression for:
        \[B_\lambda = \frac{dF}{d\lambda d\Omega}\]
        \item Starting from \(B_\lambda\), find the peak wavelength \(\lambda_{max}\) emitted by a star of temperature \(T\) (you will have to solve something numerically).
        \item At low frequencies we have the Rayleigh-Jeans approximation. Obtain an expression for \(B_\nu\) at low frequencies.
        \end{alternativas}


\end{pproblem}

\pts{3}
\begin{pproblem}
    The luminosity of a secondary body depends on the visible illuminated area of the object. In this question we briefly study that phenomenon.
    \begin{alternativas}
        \item  Consider the following situation:
        \begin{figure}[H]
            \centering
            \includegraphics[width=0.8\linewidth]{imagens/fotometria 1.png}
            \caption{Sun-Earth-Moon diagram}
        \end{figure}

        Find an expression for the ratio \(\Phi\) of the illuminated area, as a function of \(\alpha\), to the total area of the planet.
        
        \item Find the angles \(\alpha\) at which the Moon is in the phases new, waxing, full, and waning, respectively.
    \end{alternativas}


\end{pproblem}

\pts{5}
\begin{pproblem} (Magna notes)
    In this problem we model the effect of the atmosphere on Earth. Suppose that the Sun
is a blackbody of temperature \(T_1\) and radius \(R_1\). Earth is a sphere located
at a distance \(R\) from the Sun, with radius \(R_3\). Earth's emissivity is \(\epsilon_3\).
\begin{alternativas}
    \item If Earth had no atmosphere, determine its equilibrium temperature, \(T_3\).
    \item  Now we include the atmosphere. Model it as a spherical shell of gas,
    with emissivity \(\epsilon_2\) and outer radius \(R_2>R_3\), concentric with Earth. In thermal equilibrium,
    its absorptivity at ultraviolet and infrared wavelengths is \(\epsilon_2\). The
    atmosphere transmits a fraction \(t\) of the ultraviolet radiation but is completely opaque to
    infrared radiation. Assuming that the Sun emits ultraviolet light while Earth emits and re-emits
    in the infrared, determine the temperature \(T_2\) of the atmosphere and the temperature \(T_3\) of Earth in thermodynamic
    equilibrium. Assume that the atmosphere is a perfect thermal conductor, so that
    all radiation incident on it is distributed uniformly over its surface.
\end{alternativas}


\end{pproblem}


\pts{2}
\begin{pproblem}(Problem set 2 - 2021)
    The Ring Nebula (M57) has an apparent magnitude of 9 and an
angular diameter of 2$^{\prime}$ for an observer on Earth. What would be the apparent magnitude of the night
sky of a planet orbiting a star exactly at the center of M57?

\end{pproblem}


\pts{4}
\begin{pproblem} (Adapted from Problem set 4 - 2021)
    The Cherenkov effect was first detected by the Soviet scientist
    Pavel Cherenkov in 1937. Later, together with his colleagues I. E. Tamm and
    I. M. Frank, he gave the phenomenon a physical interpretation and was awarded the Nobel Prize in Physics
    in 1958. Before a mathematical treatment, we first need to understand the principle a little better.

    When high-energy charged particles travel through a dielectric medium, atoms can be excited if
    the particle speed exceeds the phase velocity ($\frac{c}{n}$). Those atoms, in
    turn, emit electromagnetic radiation when they return to the ground state. The emitted waves
    spread spherically and, when superposed, form a cone of opening angle $2\alpha$,
    as the figure below shows.
    \begin{figure}[H]
      \centering
      \includegraphics[width=0.9\linewidth]{imagens/cherenkov.png}  
      \caption{Radiation mechanism of the Cherenkov effect}
    \end{figure}

    This effect is similar to a jet moving at supersonic speed: it follows the same
    principle as the Mach cone, but with light. We now develop the mathematical model of the Cherenkov effect.

    \begin{center}
        \textbf{Part A - Theoretical Model}    
    \end{center}
    
    Consider a particle moving at relativistic speeds through a medium of refractive index
    $n$. Its rest mass is $m_0$, and it has linear momentum $p$ and speed $v$. At a
    given moment, a photon is emitted at an angle $\alpha$, as Figure 1 shows.

    \begin{alternativas}
        \item Letting $\mu$ be the frequency of the emitted photon, determine its linear
        momentum, $p_\mu$, and its energy, $E_\mu$. Your answer must be in terms of $n$, $\mu$, and physical constants.
        \item Find an expression for the linear momentum of the particle after it emits
    the photon, in terms of $p_\mu$, $p$, and $\alpha$.
        \item Letting $\beta_n = \frac{c}{vn}$, prove that the relation below holds:
        \begin{equation}
            \cos\alpha = \frac{1}{\beta_n}
        \end{equation}
        \item Assuming that linear momentum and energy are conserved, determine the
        minimum speed at which the Cherenkov effect can occur.
        \textbf{Hint:} Compared with the other parameters, the factor $(n^2-1)h\mu$ may be neglected.

    \begin{center}
        \textbf{Part B - Nuclear Reactions}    
    \end{center}
    
    The proton-proton chain is a sequence of fusion reactions that convert hydrogen into helium.
    One possible branch of the proton-proton chain is $pp \text{ IV}$, in which, in principle, a
    helium-3 atom reacts directly with a proton, according to the reaction below:

    \begin{equation}
        ^3\text{He} + ^1\text{H} \rightarrow ^4\text{He} + \nu + \_\_
    \end{equation}

        \item  State the nuclear conservation law you use, and identify the missing particle
        in the blank of the reaction above.

        \item  State the nuclear conservation law you use, and identify the missing particle
        in the blank of the reaction below:

        \begin{equation}
            \pi^- \rightarrow \mu^- + \_
        \end{equation}
        
        \textbf{Data:} Pion mass: $140 \ \text{MeV}/c^2$, muon mass: $106 \ \text{MeV}/c^2$.
    \end{alternativas}


\end{pproblem}


\pts{3} 
\begin{pproblem}(Problem set 3 - Vinhedo 2022)
Juvelino, observing directly from his observatory in Paris, France, monitors the star Polaris ($\alpha\text{UMi}$). His goal is to find the color temperature $T_c$ of the star. Some of the data he has on his target are:

\begin{itemize}
    \item Apparent magnitude in the $V$ band: $V = 1.98$;
    \item Absolute magnitude in the $V$ band: $M_V = -3.60$;
    \item Absolute magnitude in the $B$ band: $M_B = -3.19$.
\end{itemize}

Using the information provided, help Juvelino!

\begin{alternativas}
    \item From many observations, Juvelino determined that the interstellar extinction in the $V$ band toward Polaris is $a_V = 5.8 \, \text{mag/kpc}$. Find the distance, in pc, from $\alpha\text{UMi}$ to Earth.
    \item Using the empirical relation
    \begin{equation}
        \frac{A_V}{E_{B-V}} = 3.0
    \end{equation}
    where $A_V$ is the total interstellar extinction in the $V$ band and $E_{B-V}$ is the $B-V$ color excess, find the $B-V$ color index of the observed star.
    \item Show the relation
    \begin{equation}
        T_c = \frac{7009}{(B-V) + 0.47}
    \end{equation}
    in which the color temperature is given in kelvin. To do so, use the fact that stars of spectral class $A0$ have $(B-V) = 0$ and $T_c = 15000 \, \text{K}$. Also use the wavelengths of the $B$ and $V$ bands, $\lambda_B = 440 \, \text{nm}$ and $\lambda_V = 548 \, \text{nm}$, respectively. Justify any approximations you make.

    \textbf{HINT: Planck's law may be useful}
    \item Find the color temperature of Polaris.
\end{alternativas}
\end{pproblem}

\pts{4}
\begin{pproblem}
    (Problem set 4 - 2020) Frightening as it is, Juvelino's banana plantation has a perfect sky for one of his two passions: astrophotography. He has a telescope 200 mm in diameter with a super CCD attached. Its parameters are: pixel side \(5\ \mu\text{m}\), plate scale \(3.2\ \text{arcmin/mm}\), quantum efficiency 97\%, RON (read-out noise) 1 count/pixel, and DN (thermal noise) 1 count/(pixel.hour). He prefers to observe in the \(V\) band (central wavelength \(\lambda_V = 5500\ \text{Å}\), bandwidth \(\Delta \lambda = 820\ \text{Å}\)), because he remembers a magic number associated with it: the photon flux of a magnitude-0 object in this band is \(\phi_0 = 1000\ \text{photons/}(\AA\cdot \text{cm}^2)\). Juvelino has two favorite targets: a star of magnitude \(m_V = 4.51\) and a globular cluster of uniform surface brightness in \(V\) of \(19.5\ \text{mag/arcsec}^2\) and angular diameter \(\theta = 40''\).
\end{pproblem}

    \begin{alternativas}
        \item The night-sky brightness at Juvelino's farm is \(21.9\) mag/arcsec\(^2\) in the V band, and the astronomical seeing is \(\theta_0=1.2''\). Calculate the signal-to-noise ratio obtained in a photo of his favorite star with exposure time \(t_{exp}=30\)s.
        \item Calculate the zero-point magnitude of Juvelino's telescope + CCD system, using the data for the star.
        
        Juvelino's other great passion is Jiang, a Chinese university student who lives in Beijing. After years apart, Jiang suggests that Juvelino move to China. Before deciding, he looks up the night-sky brightness in the V band in his beloved's city, which is \(19.3\) mag/arcsec\(^2\). That will affect imaging of the cluster, which is also visible from Jiang's home.

        \item Juvelino cannot live without photographing the cluster. He has inviolable rules for imaging this object, wherever he does it: reach a signal-to-noise ratio $S/N=10$, and take the same number $X$ of photos of the cluster in one night. $X$ is the number of photos allowed by the sky conditions at his banana plantation. That number of frames per night is limited by the CCD battery (which lasts $3.2$ hours) and by the exposure time required for each photo. Help him with this dilemma! Do the calculations, find how many photos of the cluster he can take in one night (in Beijing and at the farm), and conclude: will he see his beloved again, or will this love be too noisy?
    \end{alternativas}


\end{document}
